\chapter*{Заключение} \label{ch-conclusion}
\addcontentsline{toc}{chapter}{Заключение}	% в оглавление

В рамках работы спроектирована, реализована и внедрена в промышленную
эксплуатацию платформы Mindbox система автоматической бесшовной ротации
криптографических ключей подписи JWT. Цель работы достигнута,
шесть задач, поставленных во введении, решены.

Анализ исходной схемы межсервисной авторизации Mindbox выявил
корневую уязвимость: единый закрытый RSA-ключ был доступен всем
160 микросервисам платформы и использовался на всех боевых контурах.
Моделирование угроз по методологии STRIDE с количественной оценкой DREAD
показало, что все семь выявленных угроз имеют высокий уровень критичности
($R_{\text{DREAD}} \geq 7{,}4$); системообразующими являются невозможность
ротации ключа~(8,2) и повышение привилегий при компрометации сервиса~(8,0).
Обзор готовых систем управления идентификацией показал, что ни одна
из рассмотренных не обеспечивает возможности встраивать дополнительную логику
в процесс ротации ключей и использовать ротацию как механизм экстренной
инвалидации токенов. Service Mesh с mTLS определён как целевое состояние
безопасности межсервисного взаимодействия, однако его внедрение требует
длительного инфраструктурного перехода, несовместимого со сроками. Это
обосновало разработку собственного решения на основе открытых стандартов
JWT, JWK, Token Exchange и OIDC Discovery.

Спроектирована и реализована система авторизации из пяти компонентов:
микросервиса auth-jwt, Kubernetes-оператора асинхронной доставки токенов
и публичных ключей, Terraform Provider, Terraform-модуля и клиентской
библиотеки Mindbox.Authentication. В сравнении с вариантом развёртывания
отдельного агента для каждого микросервиса, оператор обслуживает весь кластер
единственным экземпляром~--- расчётные требования снижаются с 145~ГБ памяти
и 46~ядер CPU до 5,5~ГБ и 2~ядер. Бесшовность ротации обеспечивает
трёхшаговый алгоритм, в котором каждый следующий шаг выполняется только
после подтверждения того, что обновлённые ключи распространились до всех
потребителей~--- это исключает переход к следующему шагу на частично
обновлённой системе. Межконтурное взаимодействие построено на Token Exchange,
конфигурация всех клиентов управляется единым Terraform-модулем,
гарантирующим однородность контуров.

Тестирование выстроено по типам компонентов: auth-jwt и клиентская
библиотека покрыты юнит- и интеграционными тестами через Testcontainers
и \texttt{Web\-Application\-Factory}; оператор~--- сквозными тестами на кластере
Kind; Terraform Provider~--- юнит-тестами маппинга. Для auth-jwt также
проведены нагрузочное тестирование в Grafana~k6 и хаос-инжиниринг
в промышленной среде. Для перевода 160 микросервисов на новую схему
авторизации разработана инкрементальная процедура из семи шагов
и реализован специализированный навык для ИИ-агента Claude Code,
автоматизирующий её выполнение; результат каждого шага проверяется
независимой языковой моделью~\cite{Zahraei2026}. Автоматизация позволила
перевести все сервисы в сжатые сроки, включая написанные на Python
и JavaScript языках, по которым у команды авторизации мало экспертизы.

Система введена в промышленную эксплуатацию 30~марта 2026~года. За два
месяца работы на всех контурах платформы не было зафиксировано ни одного
инцидента в работе сервисов~--- несмотря на то что от системы авторизации
непрерывно зависят все 160~микросервисов платформы. Автоматическая ротация
ключей выполняется по расписанию без участия инженеров. Время ответа auth-jwt
на боевом контуре Gamma по 95-му процентилю составляет 8~мс при SLO 30~мс~---
запас почти в четыре раза. Нагрузочное тестирование подтвердило пропускную
способность 15\,000~RPS при $p_{95} = 80$~мс~--- в сто раз выше промышленного
пика 150~RPS на контур. Ежедневные хаос-эксперименты с инъекцией сетевого
сбоя в зоне доступности не приводят к срабатыванию алертов SLO.
Приведённые данные подтверждают гипотезу из введения: автоматическая плановая
ротация криптографических ключей подписи без прерывания работы сервисов
и деградации производительности в платформе Mindbox достижима на основе
собственного решения, построенного на открытых стандартах и учитывающего
особенности существующей архитектуры.

Полная криптографическая изоляция контуров Beta и Gamma остаётся
открытой задачей: до её завершения Beta использует ключ Gamma.
Объективный критерий завершения~--- устойчивое нулевое значение метрики
\texttt{Contour Mismatch} (см.~\ref{subsec:ch4-rollout-isolation}).

Для команд, проектирующих собственные системы межсервисной авторизации,
из проделанной работы вытекают практические следствия. Разработку собственного
решения оправдывает лишь наличие специфических требований, не покрываемых
готовыми решениями~--- в частности, необходимость встраивать дополнительную
логику в процесс ротации ключей или использовать ротацию как механизм
экстренной инвалидации токенов; в остальных случаях рациональнее опираться
на готовые решения. При необходимости доставлять секреты во многие поды
паттерн <<Kubernetes-оператор как механизм доставки>> эффективнее
развёртывания отдельного агента на каждый сервис: в условиях Mindbox
разница составила два порядка по потреблению памяти. Пошаговый алгоритм
ротации с явными подтверждениями готовности перед каждым переходом
устойчивее фиксированного ожидания~--- ротация автоматически приостанавливается
при инфраструктурных сбоях, не допуская перехода к следующему шагу
на неполностью обновлённой системе. При Token Exchange ошибка во внешнем
контуре не должна останавливать синхронизацию текущего~--- это сохраняет
работоспособность при недоступности зависимого контура.

Среди направлений дальнейшего развития системы выделяются два приоритетных.
Первое~--- переход с RSA-SHA512 на ECDSA: при сопоставимой криптостойкости
алгоритм обеспечивает генерацию токенов в 2,8~раза быстрее и кратно меньший
размер ключей; основная сложность перехода~--- миграция всех компонентов
системы на новый алгоритм с обеспечением совместимости с внешними
потребителями. Второе~--- интеграция с автономными ИИ-агентами: по мере
их распространения в платформе потребуется выдавать им токены с чётко
ограниченными политиками доступа, управлять жизненным циклом таких токенов
и обеспечивать возможность их оперативного отзыва.
